\chapter{BF theory and discrete gauge theories}\label{ch:bf}

If Chern--Simons theory is the canonical one-field topological gauge theory in three dimensions, then BF theory is its simplest two-field companion. It is parity-even, it exists in every dimension, and in the compact abelian case it provides the cleanest continuum route to discrete gauge theory. In particular, three-dimensional compact $U(1)$ BF theory at integer level $N$ reduces to $\Z_N$ gauge theory, which is why this chapter is the natural bridge from the continuum topological actions of Chapters~\ref{ch:chern-simons} and \ref{ch:dw} to the toric-code and topological-order discussion in Chapter~\ref{ch:toporder}. We begin with the local action, its gauge invariances, and its basic observables \cite{Horowitz1989,BlauThompson1991}.

\section{The BF action, equations of motion, and observables}\label{sec:bf-action}

\begin{assumption}[Standing assumptions for Section~\ref{sec:bf-action}]\label{ass:bf-standing}
Throughout this section, $M$ is a closed oriented smooth $3$-manifold. We begin with the abelian theory, so both gauge fields are one-forms
\begin{equation*}
A,B \in \Omega^1(M).
\end{equation*}
Unless stated otherwise, the physically relevant case is the compact $U(1)$ theory, in which the holonomies of $A$ and $B$ are defined only modulo $2\pi$. Local variational formulas are first written as if $A$ and $B$ were globally defined real-valued one-forms; compactness enters when we discuss large gauge transformations and charge quantization. The boundary case is deferred, since on a manifold with boundary the action acquires extra boundary terms and edge degrees of freedom.
\end{assumption}

\begin{definition}[Abelian BF action]\label{def:bf-action}
Under Assumption~\ref{ass:bf-standing}, the three-dimensional abelian BF action at level $k$ is
\begin{equation}\label{eq:bf-action}
S_{\mathrm{BF}}[A,B]
:=
\frac{k}{2\pi}\int_M B\wedge dA.
\end{equation}
\end{definition}

\noindent The structural similarity to the abelian Chern--Simons action is immediate: the action is local and metric-independent, so it is a candidate topological field theory from the outset. The difference is that BF theory contains two gauge fields rather than one. One of them, $B$, will act as a Lagrange multiplier enforcing the flatness of $A$; by symmetry of the variational problem, $A$ will in turn force $B$ to be closed.

\medskip

\noindent If $A$ and $B$ are treated as noncompact $\R$-valued one-forms, then \eqref{eq:bf-action} makes sense for any real $k$. In the compact $U(1)$ theory relevant for this review, the path-integral phase $e^{iS_{\mathrm{BF}}}$ must be well defined under large gauge transformations, and this requires
\begin{equation}\label{eq:bf-level-integral}
k\in \Z.
\end{equation}
In Sections~6.2--6.4 we will rename this integer level by $N$ to emphasize its interpretation as the order of the residual discrete gauge group.

\begin{proposition}[Gauge invariance]\label{prop:bf-gauge}
For the abelian theory on a closed manifold, the action \eqref{eq:bf-action} is invariant under the small gauge transformations
\begin{equation}\label{eq:bf-gauge-transformations}
A \longmapsto A + d\lambda,
\qquad
B \longmapsto B + d\chi,
\end{equation}
where $\lambda,\chi \in \Omega^0(M)$.
\end{proposition}

\begin{proof}
Substituting \eqref{eq:bf-gauge-transformations} into \eqref{eq:bf-action} gives
\begin{align*}
S_{\mathrm{BF}}[A+d\lambda,B+d\chi]
&=
\frac{k}{2\pi}\int_M (B+d\chi)\wedge d(A+d\lambda) \\
&=
\frac{k}{2\pi}\int_M (B+d\chi)\wedge dA,
\end{align*}
since $d^2\lambda=0$. Expanding the remaining term yields
\begin{equation*}
S_{\mathrm{BF}}[A+d\lambda,B+d\chi]
=
S_{\mathrm{BF}}[A,B]
+
\frac{k}{2\pi}\int_M d\chi\wedge dA.
\end{equation*}
Now
\begin{equation*}
d(\chi\, dA)=d\chi\wedge dA + \chi\, d^2A = d\chi\wedge dA,
\end{equation*}
so the last term is an integral of an exact $3$-form. Because $M$ is closed, Stokes' theorem gives
\begin{equation*}
\int_M d\chi\wedge dA = \int_M d(\chi\, dA)=0.
\end{equation*}
Hence
\begin{equation*}
S_{\mathrm{BF}}[A+d\lambda,B+d\chi]=S_{\mathrm{BF}}[A,B].
\end{equation*}
\end{proof}

\begin{remark}[Compact versus noncompact theories]\label{rem:bf-compact}
For noncompact $\R$-valued gauge fields, Proposition~\ref{prop:bf-gauge} is the whole story: $k$ may be any real number, and charges need not be quantized. In the compact $U(1)$ theory, the holonomies of $A$ and $B$ are angular variables, so large gauge transformations can shift them by integer multiples of $2\pi$. Demanding that the exponentiated action and the exponentiated observables below be single-valued is what forces the integrality conditions
\begin{equation*}
k,q,n \in \Z.
\end{equation*}
Since the later reduction to $\Z_N$ gauge theory uses exactly this compact case, all subsequent BF sections work in that regime.
\end{remark}

\begin{proposition}[Equations of motion]\label{prop:bf-eom}
Under Assumption~\ref{ass:bf-standing}, the Euler--Lagrange equations of the abelian BF action are
\begin{equation}\label{eq:bf-eom}
dA = 0,
\qquad
dB = 0.
\end{equation}
\end{proposition}

\begin{proof}
Varying \eqref{eq:bf-action} gives
\begin{align*}
\delta S_{\mathrm{BF}}
&=
\frac{k}{2\pi}\int_M \left(\delta B\wedge dA + B\wedge d(\delta A)\right).
\end{align*}
To treat the second term, use the Leibniz rule with $\deg(B)=1$:
\begin{equation*}
d(B\wedge \delta A)=dB\wedge \delta A - B\wedge d(\delta A).
\end{equation*}
Rearranging,
\begin{equation*}
B\wedge d(\delta A)=dB\wedge \delta A - d(B\wedge \delta A).
\end{equation*}
Substituting this into $\delta S_{\mathrm{BF}}$ yields
\begin{align*}
\delta S_{\mathrm{BF}}
&=
\frac{k}{2\pi}\int_M \left(\delta B\wedge dA + dB\wedge \delta A - d(B\wedge \delta A)\right).
\end{align*}
The last term integrates to zero because $M$ is closed. Since $dB$ is a two-form and $\delta A$ is a one-form, we have
\begin{equation*}
dB\wedge \delta A = \delta A \wedge dB.
\end{equation*}
Therefore
\begin{equation}\label{eq:bf-variation}
\delta S_{\mathrm{BF}}
=
\frac{k}{2\pi}\int_M \left(\delta B\wedge dA + \delta A\wedge dB\right).
\end{equation}
The variations $\delta A$ and $\delta B$ are independent, so stationarity of the action implies
\begin{equation*}
dA=0,
\qquad
dB=0.
\end{equation*}
\end{proof}

\noindent Proposition~\ref{prop:bf-eom} shows that the classical solutions are pairs of flat abelian gauge fields modulo gauge equivalence. As in Chern--Simons theory, there are no local propagating bulk degrees of freedom; the physically relevant information is global holonomy data.

\subsection{Observables}

The natural observables are easiest to state first in arbitrary dimension. In $d$ spacetime dimensions, abelian BF theory contains a one-form $A$ and a $(d-2)$-form $B$, so the basic exponentiated defects are
\begin{equation}\label{eq:bf-general-observables}
W_q(C) := \exp\!\left(iq\int_C A\right),
\qquad
T_n(\Sigma^{d-2}) := \exp\!\left(in\int_{\Sigma^{d-2}} B\right),
\end{equation}
where $C$ is a closed oriented one-cycle and $\Sigma^{d-2}$ is a closed oriented $(d-2)$-cycle. In three dimensions, $d-2=1$, so the second observable is \emph{also} supported on a closed curve. We therefore specialize \eqref{eq:bf-general-observables} to
\begin{equation}\label{eq:bf-line-operators}
W_q(C) := \exp\!\left(iq\oint_C A\right),
\qquad
T_n(C') := \exp\!\left(in\oint_{C'} B\right),
\end{equation}
with $q,n\in \Z$ in the compact theory.

\begin{proposition}[Gauge invariance of the line operators]\label{prop:bf-obs-gauge}
If $C$ and $C'$ are closed oriented curves in $M$, then the operators \eqref{eq:bf-line-operators} are invariant under the small gauge transformations \eqref{eq:bf-gauge-transformations}.
\end{proposition}

\begin{proof}
Under $A\mapsto A+d\lambda$,
\begin{equation*}
\oint_C (A+d\lambda)
=
\oint_C A + \oint_C d\lambda
=
\oint_C A,
\end{equation*}
because the integral of an exact form over a closed curve vanishes. Hence $W_q(C)$ is gauge invariant. The same argument with $B\mapsto B+d\chi$ shows that
\begin{equation*}
\oint_{C'} (B+d\chi)=\oint_{C'}B,
\end{equation*}
so $T_n(C')$ is gauge invariant as well.
\end{proof}

\noindent The mixed correlator of these two operators is the BF analogue of the abelian Wilson-loop calculation in Chapter~\ref{ch:chern-simons}. Section~6.3 will show that, in the compact level-$N$ theory, their expectation value is a pure linking-number phase. That result is the direct continuum precursor of the mutual braiding phase in $\Z_N$ gauge theory and, for $N=2$, of the electric--magnetic braiding in the toric code.

\begin{remark}[Nonabelian BF theory]\label{rem:nonabelian-bf}
There is a nonabelian generalization in which $A$ is a connection on a principal $G$-bundle, $B$ is an adjoint-valued one-form, and the action is
\begin{equation*}
S_{\mathrm{BF}}^{\mathrm{nonab}}
=
\frac{k}{2\pi}\int_M \tr(B\wedge F_A),
\qquad
F_A = dA + A\wedge A.
\end{equation*}
Its equations of motion are
\begin{equation*}
F_A=0,
\qquad
d_A B = 0.
\end{equation*}
We record this only for orientation. The rest of this chapter stays in the compact abelian theory, because that is the case directly related to $\Z_N$ gauge theory and to the applications developed later in the paper.
\end{remark}

\section{Compact BF at level \texorpdfstring{$N$}{N} and the emergence of \texorpdfstring{$\Z_N$}{Z\_N} gauge theory}\label{sec:bf-zn}

We now specialize to the compact theory at integer level
\begin{equation*}
k=N\in \Z_{>0}.
\end{equation*}
The claim of this section is that compact $U(1)$ BF theory at level $N$ does not describe a continuous infrared gauge group. Instead, after the periodicity of the fields is properly taken into account, only $N$-torsion holonomies survive. That is exactly the field content of $\Z_N$ gauge theory.

\subsection{The finite-dimensional Fourier identity}

The mechanism is the same one that relates a compact angular variable to an integer Fourier mode. The one-dimensional identity we need is
\begin{equation}\label{eq:periodic-delta}
\sum_{m\in \Z} e^{iNm\theta}
=
\frac{2\pi}{N}\sum_{r=0}^{N-1}\sum_{\ell\in \Z}
\delta\!\left(\theta-\frac{2\pi r}{N}-2\pi \ell\right).
\end{equation}
This is just Poisson resummation for a periodic function. It says that summing over all integer Fourier modes of period $N$ forces the angle $\theta$ to lie at one of the $N$ roots of unity.

\medskip

\noindent Equation \eqref{eq:periodic-delta} is the precise reason a compact BF path integral produces discrete gauge data: once the compact field is integrated or summed over, the conjugate holonomy variable is projected onto $N$-th roots of unity.

\subsection{Lattice regularization of the compact theory}

To make that statement explicit, choose a finite oriented cell decomposition $K$ of the closed $3$-manifold $M$. We place a compact $U(1)$ variable
\begin{equation*}
U_e = e^{iA_e}\in U(1)
\end{equation*}
on each oriented $1$-cell $e\in K^1$. For every oriented $2$-cell $p\in K^2$, define the plaquette holonomy
\begin{equation}\label{eq:plaquette-holonomy}
U_p := \prod_{e\subset \partial p} U_e^{\epsilon_{p,e}}
=
\exp\!\left(i\sum_{e\subset \partial p}\epsilon_{p,e}A_e\right),
\end{equation}
where $\epsilon_{p,e}=\pm 1$ is the incidence sign of $e$ in the boundary of $p$.

\medskip

\noindent The compact BF field $B$ is the variable Fourier-dual to the plaquette holonomy. A convenient lattice realization is therefore to write the partition sum as
\begin{equation}\label{eq:lattice-bf-partition}
Z_{\mathrm{BF},N}(K)
:=
\int\!\prod_{e\in K^1} dU_e
\prod_{p\in K^2}
\left(\sum_{m_p\in \Z} U_p^{\,Nm_p}\right),
\end{equation}
where $dU_e$ is the normalized Haar measure on $U(1)$ and the integer $m_p$ is the Fourier mode conjugate to the compact $B$ variable on the cell dual to $p$. Equation \eqref{eq:lattice-bf-partition} is the lattice form of the compact path integral: the sum over all integers $m_p$ is the discrete Fourier expansion of the periodic delta functional implementing the compact $B$ constraint.

\begin{proposition}[Projection to $N$-torsion plaquette holonomy]\label{prop:bf-root-unity}
For each plaquette $p$, the factor
\begin{equation}\label{eq:root-unity-projector}
\sum_{m_p\in \Z} U_p^{\,Nm_p}
\end{equation}
vanishes unless $U_p$ is an $N$-th root of unity. Equivalently, the lattice BF partition function localizes on configurations satisfying
\begin{equation}\label{eq:Up-N}
U_p^N = 1
\qquad
\text{for every plaquette }p.
\end{equation}
\end{proposition}

\begin{proof}
Write
\begin{equation*}
U_p = e^{i\theta_p},
\qquad
\theta_p \in \R/2\pi\Z.
\end{equation*}
Then
\begin{equation*}
U_p^{\,Nm_p}=e^{iNm_p\theta_p},
\end{equation*}
so the plaquette factor in \eqref{eq:lattice-bf-partition} is exactly the left-hand side of \eqref{eq:periodic-delta} with $\theta=\theta_p$. Therefore it is supported only when
\begin{equation*}
\theta_p = \frac{2\pi r_p}{N}
\qquad
\text{for some }r_p\in\{0,1,\dots,N-1\},
\end{equation*}
which is equivalent to
\begin{equation*}
U_p = e^{2\pi i r_p/N}.
\end{equation*}
This is precisely the condition \eqref{eq:Up-N}.
\end{proof}

\noindent Proposition~\ref{prop:bf-root-unity} shows the essential kinematical point: compactness reduces the continuous plaquette variable to $N$ discrete possibilities. The residual gauge data therefore no longer takes values in all of $U(1)$, but only in the subgroup
\begin{equation}\label{eq:muN}
\mu_N := \left\{e^{2\pi i r/N}\mid r=0,1,\dots,N-1\right\}\cong \Z_N.
\end{equation}

\subsection{From \texorpdfstring{$\mu_N$}{mu\_N} holonomies to \texorpdfstring{$\Z_N$}{Z\_N} gauge fields}

Now combine this with Proposition~\ref{prop:bf-eom}, which already showed that the BF equations of motion impose local flatness. On a simply connected patch, a flat connection is pure gauge, so after quotienting by the remaining continuous gauge redundancy one may choose representatives whose transition functions take values in $\mu_N$. Equivalently, one can describe the residual field by a $\Z_N$-valued lattice $1$-cochain.

\medskip

\noindent Concretely, choose, for each edge $e$, an integer label
\begin{equation*}
h_e \in \Z_N
\end{equation*}
and write
\begin{equation}\label{eq:U-from-h}
U_e = e^{2\pi i h_e/N}.
\end{equation}
Then the plaquette holonomy \eqref{eq:plaquette-holonomy} becomes
\begin{equation}\label{eq:Up-from-h}
U_p
=
\exp\!\left(\frac{2\pi i}{N}\sum_{e\subset \partial p}\epsilon_{p,e}h_e\right)
=
\exp\!\left(\frac{2\pi i}{N}(dh)_p\right),
\end{equation}
where $h$ is now viewed as a $\Z_N$-valued $1$-cochain and $d$ is the lattice coboundary operator. Thus
\begin{equation}\label{eq:flatness-mod-N}
U_p = 1
\qquad\Longleftrightarrow\qquad
(dh)_p=0 \ \text{in }\Z_N.
\end{equation}
In other words, a flat $\Z_N$ lattice gauge field is exactly a compact $U(1)$ field whose holonomies lie in $\mu_N$ and whose plaquette curvature is trivial in that subgroup.

\medskip

\noindent This is the promised reduction: the compact BF path integral does not retain arbitrary flat $U(1)$ connections, but only their $\mu_N\cong \Z_N$ reduction. In other words, the surviving configurations are flat $\Z_N$ connections. That is the untwisted discrete gauge theory that Chapter~\ref{ch:dw} will later reinterpret as Dijkgraaf--Witten theory with trivial cocycle, and in the special $\Z_2$ case Section~\ref{sec:dw-toric} will match to the toric-code phase.

\begin{corollary}[Charges are defined modulo \texorpdfstring{$N$}{N}]\label{cor:charges-mod-N}
In compact level-$N$ BF theory, the line operators of Section~\ref{sec:bf-action} depend only on the charges modulo $N$:
\begin{equation}\label{eq:charges-mod-N}
W_{q+N}(C)=W_q(C),
\qquad
T_{n+N}(C')=T_n(C').
\end{equation}
\end{corollary}

\begin{proof}
After the reduction to $\Z_N$ gauge data, the line operator $W_q(C)$ can be written in terms of the $\Z_N$-valued edge labels introduced above:
\begin{equation*}
\!\!\!
W_q(C)
=
\exp\!\left(\frac{2\pi i q}{N}\sum_{e\subset C}\epsilon_{C,e}h_e\right).
\end{equation*}
The exponent changes by an integer multiple of $2\pi i$ when $q$ is shifted by $N$, so $W_{q+N}(C)=W_q(C)$. The same argument applied to the dual $\Z_N$ data for $B$ gives $T_{n+N}(C')=T_n(C')$.
\end{proof}

\begin{remark}[What is proved here and what is postponed]\label{rem:bf-zn-scope}
Section~\ref{sec:bf-zn} proves the structural statement needed for the rest of the paper: compact $U(1)$ BF theory at level $N$ reduces the continuous gauge data to flat $\Z_N$ gauge data, and the defect charges are correspondingly defined only modulo $N$. Two more concrete consequences are deferred to the next sections:
\begin{itemize}
\item Section~6.3 computes the mixed Wilson/BF correlator and shows that it produces the expected $\Z_N$ linking phase.
\item Section~6.4 quantizes the reduced phase space on $\Sigma_g$ and derives the genus-$g$ ground-state degeneracy $N^{2g}$.
\item Appendix~\ref{app:bf-zn}, especially Section~\ref{sec:appD-bf-lattice}, will later supply the full lattice bookkeeping behind the compact-BF $\leftrightarrow \Z_N$ equivalence and the direct comparison with the untwisted Dijkgraaf--Witten sum.
\end{itemize}
Those two results are the cleanest checks that the compact BF theory and the untwisted $\Z_N$ gauge theory are not merely morally related, but actually have the same topological content \cite{Horowitz1989,GKSW2015,Bhardwaj2023Lectures}.
\end{remark}

\section{The mixed BF linking phase}\label{sec:bf-linking}

The analogue of the abelian Chern--Simons Wilson-loop calculation from Example~\ref{ex:abelian-wilson} is now mixed rather than diagonal: one line operator is built from $A$, the other from $B$. In three dimensions these are both supported on closed curves, and their correlator detects the linking number.

\begin{assumption}[Setup for the explicit BF linking calculation]\label{ass:bf-linking}
For the remainder of this section, take $M=S^3$ and let $C,C'\subset S^3$ be disjoint oriented closed curves with charges $q,n\in \Z$. Because every knot in $S^3$ is null-homologous, we may choose Seifert surfaces $\Sigma$ and $\Sigma'$ satisfying
\begin{equation*}
\partial \Sigma = C,
\qquad
\partial \Sigma' = C'.
\end{equation*}
Let $\delta_C,\delta_{C'}$ denote the associated distributional two-form currents, and let
\begin{equation*}
\eta_C:=\delta_\Sigma,
\qquad
\eta_{C'}:=\delta_{\Sigma'}
\end{equation*}
be the corresponding one-form currents, so that
\begin{equation}\label{eq:bf-current-boundaries}
d\eta_C=\delta_C,
\qquad
d\eta_{C'}=\delta_{C'}.
\end{equation}
\end{assumption}

\begin{proposition}[Mixed linking phase in compact abelian BF theory]\label{prop:bf-linking}
Under Assumption~\ref{ass:bf-linking}, the normalized correlator of the two basic line operators in compact level-$N$ BF theory is
\begin{equation}\label{eq:bf-linking-phase}
\left\langle W_q(C)\,T_n(C')\right\rangle_{S^3}
=
\exp\!\left(-\frac{2\pi i q n}{N}\,\mathrm{Lk}(C,C')\right),
\end{equation}
with the sign convention determined by the BF action \eqref{eq:bf-action}. Reversing the orientation of $S^3$ or replacing $S_{\mathrm{BF}}$ by $-S_{\mathrm{BF}}$ complex-conjugates the phase.
\end{proposition}

\begin{proof}
By definition,
\begin{equation}\label{eq:bf-correlator-start}
\left\langle W_q(C)\,T_n(C')\right\rangle_{S^3}
:=
\frac{1}{Z_{\mathrm{BF},N}(S^3)}
\int \mathcal{D}A\,\mathcal{D}B\;
\exp\!\left(
i\frac{N}{2\pi}\int_{S^3} B\wedge dA
+
iq\oint_C A
+
in\oint_{C'} B
\right).
\end{equation}
Using the distributional currents from Assumption~\ref{ass:bf-linking}, rewrite the line integrals as
\begin{equation*}
\oint_C A = \int_{S^3} A\wedge \delta_C,
\qquad
\oint_{C'} B = \int_{S^3} B\wedge \delta_{C'}.
\end{equation*}
Hence
\begin{equation}\label{eq:bf-correlator-currents}
\left\langle W_q(C)\,T_n(C')\right\rangle_{S^3}
=
\frac{1}{Z_{\mathrm{BF},N}(S^3)}
\int \mathcal{D}A\,\mathcal{D}B\;
\exp\!\left(
i\frac{N}{2\pi}\int_{S^3} B\wedge dA
+
iq\int_{S^3} A\wedge \delta_C
+
in\int_{S^3} B\wedge \delta_{C'}
\right).
\end{equation}

\medskip

\noindent We now absorb the $B$-source by shifting $A$:
\begin{equation}\label{eq:bf-A-shift}
\widetilde{A}
:=
A+\frac{2\pi n}{N}\eta_{C'}.
\end{equation}
Using \eqref{eq:bf-current-boundaries},
\begin{equation*}
d\widetilde{A}
=
dA+\frac{2\pi n}{N}\delta_{C'}.
\end{equation*}
Therefore
\begin{align}
\frac{N}{2\pi}\int_{S^3} B\wedge d\widetilde{A}
&=
\frac{N}{2\pi}\int_{S^3} B\wedge dA
+
n\int_{S^3} B\wedge \delta_{C'}.
\label{eq:bf-shifted-kinetic}
\end{align}
Substituting \eqref{eq:bf-A-shift} and \eqref{eq:bf-shifted-kinetic} into \eqref{eq:bf-correlator-currents} gives
\begin{align}
\left\langle W_q(C)\,T_n(C')\right\rangle_{S^3}
&=
\frac{1}{Z_{\mathrm{BF},N}(S^3)}
\int \mathcal{D}\widetilde{A}\,\mathcal{D}B\;
\exp\!\left(
i\frac{N}{2\pi}\int_{S^3} B\wedge d\widetilde{A}
+
iq\int_{S^3} \widetilde{A}\wedge \delta_C
-i\frac{2\pi q n}{N}\int_{S^3} \eta_{C'}\wedge \delta_C
\right) \notag\\
&=
\exp\!\left(
-i\frac{2\pi q n}{N}\int_{S^3} \eta_{C'}\wedge \delta_C
\right)
\frac{1}{Z_{\mathrm{BF},N}(S^3)}
\int \mathcal{D}\widetilde{A}\,\mathcal{D}B\;
\exp\!\left(
i\frac{N}{2\pi}\int_{S^3} B\wedge d\widetilde{A}
+
iq\int_{S^3} \widetilde{A}\wedge \delta_C
\right).
\label{eq:bf-correlator-after-shift}
\end{align}

\medskip

\noindent We now evaluate the remaining functional integral. Since $B$ appears only linearly, integrating over $B$ imposes the delta-functional constraint
\begin{equation}\label{eq:bf-delta-constraint}
d\widetilde{A}=0.
\end{equation}
On $S^3$ every flat $U(1)$ connection is pure gauge because
\begin{equation*}
H^1(S^3;\R)=0.
\end{equation*}
Hence any field satisfying \eqref{eq:bf-delta-constraint} may be written as
\begin{equation*}
\widetilde{A}=d\lambda,
\end{equation*}
and therefore
\begin{equation*}
\int_{S^3}\widetilde{A}\wedge \delta_C
=
\oint_C d\lambda
=
0.
\end{equation*}
The remaining integral in \eqref{eq:bf-correlator-after-shift} is thus exactly the vacuum partition function $Z_{\mathrm{BF},N}(S^3)$, so the ratio collapses to the phase factor alone:
\begin{equation}\label{eq:bf-correlator-phase-only}
\left\langle W_q(C)\,T_n(C')\right\rangle_{S^3}
=
\exp\!\left(
-i\frac{2\pi q n}{N}\int_{S^3} \eta_{C'}\wedge \delta_C
\right).
\end{equation}

\medskip

\noindent It remains only to identify the current integral with the linking number. Since $\eta_{C'}=\delta_{\Sigma'}$ is the current associated to a Seifert surface for $C'$, we have
\begin{equation*}
\int_{S^3}\eta_{C'}\wedge \delta_C
=
\int_{\Sigma'} \delta_C
=
\Sigma'\cdot C
=
\mathrm{Lk}(C,C').
\end{equation*}
Substituting this into \eqref{eq:bf-correlator-phase-only} proves \eqref{eq:bf-linking-phase}.
\end{proof}

\begin{remark}[Relation to Chern--Simons and to the toric code]\label{rem:bf-linking-physics}
Proposition~\ref{prop:bf-linking} is the BF analogue of Example~\ref{ex:abelian-wilson} from the Chern--Simons chapter. The Chern--Simons correlator links two Wilson lines built from the same gauge field; the BF correlator links an $A$-line with a $B$-line. In both cases the answer is a pure topological phase determined by the linking number.

\medskip

\noindent The special case
\begin{equation*}
N=2,
\qquad
q=n=1
\end{equation*}
gives the phase
\begin{equation*}
\left\langle W_1(C)\,T_1(C')\right\rangle_{S^3}
=
(-1)^{\mathrm{Lk}(C,C')}.
\end{equation*}
This is the continuum precursor of the mutual minus sign acquired when the electric and magnetic quasiparticles of the toric code braid around one another. Chapter~\ref{ch:toporder} will return to exactly that phase from the lattice-model side.
\end{remark}

\section{Ground-state degeneracy on \texorpdfstring{$\Sigma_g$}{Sigma\_g}}\label{sec:bf-gsd}

The last basic check is the Hilbert-space dimension on a closed spatial surface of genus $g$. In Chern--Simons theory, the abelian torus Hilbert space had dimension $k$ because one pair of canonically conjugate holonomies survived the flatness constraint. In BF theory there are \emph{two} independent flat gauge fields, so each handle contributes two such pairs, and the Hilbert-space dimension grows as $N^{2g}$ rather than $N^g$.

\begin{assumption}[Setup for canonical quantization]\label{ass:bf-gsd}
Let
\begin{equation*}
M=\R_t\times \Sigma_g
\end{equation*}
with $\Sigma_g$ a closed oriented surface of genus $g$. We work in compact abelian BF theory at level $N\in \Z_{>0}$. Choose a symplectic basis of one-cycles
\begin{equation*}
a_1,\dots,a_g,\; b_1,\dots,b_g
\end{equation*}
for $H_1(\Sigma_g,\Z)$, and choose closed one-forms
\begin{equation*}
\alpha_1,\dots,\alpha_g,\; \beta_1,\dots,\beta_g
\end{equation*}
representing the dual cohomology basis, normalized by
\begin{equation}\label{eq:bf-basis-periods}
\oint_{a_j}\alpha_i=\delta_{ij},
\qquad
\oint_{b_j}\beta_i=\delta_{ij},
\qquad
\oint_{a_j}\beta_i=\oint_{b_j}\alpha_i=0,
\end{equation}
and
\begin{equation}\label{eq:bf-basis-pairing}
\int_{\Sigma_g}\alpha_i\wedge \beta_j=\delta_{ij},
\qquad
\int_{\Sigma_g}\beta_i\wedge \alpha_j=-\delta_{ij}.
\end{equation}
\end{assumption}

\subsection{Hamiltonian form of the BF action}

Write
\begin{equation*}
A=A_0\,dt+A_\Sigma,
\qquad
B=B_0\,dt+B_\Sigma,
\end{equation*}
and decompose the exterior derivative as
\begin{equation*}
d=dt\,\del_t+d_\Sigma.
\end{equation*}
Then
\begin{equation*}
dA=dt\wedge(\del_t A_\Sigma-d_\Sigma A_0)+d_\Sigma A_\Sigma.
\end{equation*}
Substituting this into the BF action gives
\begin{align}
B\wedge dA
&=
(B_0dt+B_\Sigma)\wedge\left(dt\wedge(\del_t A_\Sigma-d_\Sigma A_0)+d_\Sigma A_\Sigma\right) \notag\\
&=
dt\wedge\left(B_0\,d_\Sigma A_\Sigma-B_\Sigma\wedge \del_t A_\Sigma+B_\Sigma\wedge d_\Sigma A_0\right),
\label{eq:bf-split-first}
\end{align}
because $dt\wedge dt=0$ and $B_\Sigma\wedge d_\Sigma A_\Sigma$ vanishes identically on the two-dimensional surface $\Sigma_g$.

\medskip

\noindent The last term may be reorganized using
\begin{equation*}
d_\Sigma(B_\Sigma A_0)=d_\Sigma B_\Sigma\,A_0-B_\Sigma\wedge d_\Sigma A_0.
\end{equation*}
Since $A_0$ is a zero-form,
\begin{equation*}
B_\Sigma\wedge d_\Sigma A_0=A_0\, d_\Sigma B_\Sigma-d_\Sigma(B_\Sigma A_0).
\end{equation*}
On the closed surface $\Sigma_g$, the total derivative integrates to zero, so the action becomes
\begin{equation}\label{eq:bf-hamiltonian}
S_{\mathrm{BF}}
=
\frac{N}{2\pi}\int_\R dt\int_{\Sigma_g}
\left(
-B_\Sigma\wedge \del_t A_\Sigma
+B_0\,d_\Sigma A_\Sigma
+A_0\,d_\Sigma B_\Sigma
\right).
\end{equation}

\noindent Equation \eqref{eq:bf-hamiltonian} makes the constraint structure manifest. The fields $A_0$ and $B_0$ appear without time derivatives, so they are Lagrange multipliers enforcing
\begin{equation}\label{eq:bf-flat-constraints}
d_\Sigma A_\Sigma=0,
\qquad
d_\Sigma B_\Sigma=0.
\end{equation}
Thus the reduced phase space is the space of pairs of flat $U(1)$ connections modulo gauge transformations.

\subsection{Holonomy coordinates on the reduced phase space}

Every flat abelian connection is gauge-equivalent to a harmonic representative determined by its periods around the $a_i$ and $b_i$ cycles. Under Assumption~\ref{ass:bf-gsd}, we may therefore write
\begin{equation}\label{eq:bf-flat-expansion}
A_\Sigma=\sum_{i=1}^g \bigl(u_i(t)\alpha_i+v_i(t)\beta_i\bigr),
\qquad
B_\Sigma=\sum_{i=1}^g \bigl(x_i(t)\alpha_i+y_i(t)\beta_i\bigr),
\end{equation}
where
\begin{equation}\label{eq:bf-holonomies}
u_i=\oint_{a_i}A,
\quad
v_i=\oint_{b_i}A,
\qquad
x_i=\oint_{a_i}B,
\quad
y_i=\oint_{b_i}B.
\end{equation}
Because $A$ and $B$ are compact $U(1)$ gauge fields, each of these holonomies is an angular variable:
\begin{equation}\label{eq:bf-periodic-holonomies}
u_i,v_i,x_i,y_i \in \R/2\pi\Z.
\end{equation}

\medskip

\noindent Substituting \eqref{eq:bf-flat-expansion} into the kinetic term of \eqref{eq:bf-hamiltonian} gives
\begin{align*}
\int_{\Sigma_g} B_\Sigma\wedge \del_t A_\Sigma
&=
\sum_{i,j=1}^g
\int_{\Sigma_g}
\bigl(x_i\alpha_i+y_i\beta_i\bigr)\wedge
\bigl(\dot u_j\alpha_j+\dot v_j\beta_j\bigr) \\
&=
\sum_{i,j=1}^g
\left(
x_i\dot v_j\int_{\Sigma_g}\alpha_i\wedge\beta_j
+
y_i\dot u_j\int_{\Sigma_g}\beta_i\wedge\alpha_j
\right) \\
&=
\sum_{i=1}^g \left(x_i\dot v_i-y_i\dot u_i\right).
\end{align*}
Hence the reduced action is
\begin{equation}\label{eq:bf-reduced-action}
S_{\mathrm{red}}
=
\frac{N}{2\pi}\int_\R dt\sum_{i=1}^g \left(y_i\dot u_i-x_i\dot v_i\right).
\end{equation}

\noindent We may therefore read off the canonical momenta:
\begin{equation}\label{eq:bf-canonical-momenta}
p_{u_i}=\frac{N}{2\pi}y_i,
\qquad
p_{v_i}=-\frac{N}{2\pi}x_i.
\end{equation}
The key difference from Chern--Simons is now visible: each handle contributes \emph{two} canonical pairs,
\begin{equation*}
(u_i,y_i)
\qquad\text{and}\qquad
(v_i,-x_i).
\end{equation*}

\subsection{Quantization and the finite holonomy algebra}

Choose the polarization in which wavefunctions depend on the angular variables
\begin{equation*}
u_1,\dots,u_g,v_1,\dots,v_g.
\end{equation*}
Periodicity implies
\begin{equation*}
\psi(\dots,u_i+2\pi,\dots,v_j,\dots)=\psi(\dots,u_i,\dots,v_j,\dots),
\qquad
\psi(\dots,u_i,\dots,v_j+2\pi,\dots)=\psi(\dots,u_i,\dots,v_j,\dots).
\end{equation*}
Thus the momentum operators
\begin{equation*}
\hat p_{u_i}=-i\frac{\del}{\del u_i},
\qquad
\hat p_{v_i}=-i\frac{\del}{\del v_i}
\end{equation*}
have integer eigenvalues, with Fourier basis
\begin{equation}\label{eq:bf-fourier-basis}
\psi_{\vec r,\vec s}(u,v)
:=
\exp\!\left(i\sum_{i=1}^g r_i u_i + i\sum_{i=1}^g s_i v_i\right),
\qquad
r_i,s_i\in \Z.
\end{equation}

\medskip

\noindent At the same time, \eqref{eq:bf-canonical-momenta} and the periodicity \eqref{eq:bf-periodic-holonomies} imply that the classical momenta are defined only modulo $N$:
\begin{equation}\label{eq:bf-momentum-mod-N}
p_{u_i}\sim p_{u_i}+N,
\qquad
p_{v_i}\sim p_{v_i}+N.
\end{equation}
Therefore the quantum momentum labels are identified by
\begin{equation}\label{eq:bf-label-identification}
r_i\sim r_i+N,
\qquad
s_i\sim s_i+N.
\end{equation}
One convenient set of representatives is
\begin{equation}\label{eq:bf-basis-states}
r_i,s_i\in\{0,1,\dots,N-1\}.
\end{equation}
The inequivalent physical states are therefore labeled by two independent $g$-tuples in $\Z_N^g$.

\medskip

\noindent This count may be packaged as a pair of commuting finite Heisenberg algebras. Define, for each $i$,
\begin{equation}\label{eq:bf-heisenberg-ops}
\hat U_i:=e^{iu_i},
\qquad
\hat X_i:=e^{2\pi i \hat p_{u_i}/N},
\qquad
\hat V_i:=e^{iv_i},
\qquad
\hat Y_i:=e^{2\pi i \hat p_{v_i}/N}.
\end{equation}
Acting on the basis \eqref{eq:bf-fourier-basis}, these satisfy
\begin{equation*}
\hat U_i\psi_{\vec r,\vec s}=\psi_{\vec r+\mathbf e_i,\vec s},
\qquad
\hat X_i\psi_{\vec r,\vec s}=e^{2\pi i r_i/N}\psi_{\vec r,\vec s},
\end{equation*}
and
\begin{equation*}
\hat V_i\psi_{\vec r,\vec s}=\psi_{\vec r,\vec s+\mathbf e_i},
\qquad
\hat Y_i\psi_{\vec r,\vec s}=e^{2\pi i s_i/N}\psi_{\vec r,\vec s},
\end{equation*}
where $\mathbf e_i$ is the $i$th standard basis vector and all additions are understood modulo $N$. Hence
\begin{equation}\label{eq:bf-heisenberg-relations}
\hat U_i\hat X_j=e^{-2\pi i\delta_{ij}/N}\hat X_j\hat U_i,
\qquad
\hat V_i\hat Y_j=e^{-2\pi i\delta_{ij}/N}\hat Y_j\hat V_i,
\end{equation}
while all mixed pairs commute. Each pair gives the standard $N$-dimensional irreducible representation of a finite Heisenberg algebra.

\begin{proposition}[BF ground-state degeneracy on \texorpdfstring{$\Sigma_g$}{Sigma\_g}]\label{prop:bf-gsd}
For compact abelian BF theory at level $N$ on a closed oriented surface $\Sigma_g$ of genus $g$,
\begin{equation}\label{eq:bf-gsd}
\dim \mathcal{H}_{\mathrm{BF},N}(\Sigma_g)=N^{2g}.
\end{equation}
\end{proposition}

\begin{proof}
Equation \eqref{eq:bf-basis-states} shows that the physical basis is labeled by
\begin{equation*}
(\vec r,\vec s)\in \Z_N^g\times \Z_N^g.
\end{equation*}
There are $N^g$ choices for $\vec r$ and independently $N^g$ choices for $\vec s$, so the number of inequivalent states is
\begin{equation*}
N^g\cdot N^g=N^{2g}.
\end{equation*}
\end{proof}

\begin{remark}[Sanity checks and forward pointers]\label{rem:bf-gsd-checks}
The formula \eqref{eq:bf-gsd} has the expected limiting cases:
\begin{equation*}
g=0
\quad\Longrightarrow\quad
\dim \mathcal{H}_{\mathrm{BF},N}(S^2)=1,
\end{equation*}
since there are no nontrivial one-cycles on the sphere. For the torus,
\begin{equation*}
g=1
\quad\Longrightarrow\quad
\dim \mathcal{H}_{\mathrm{BF},N}(T^2)=N^2.
\end{equation*}
In particular, for the $\Z_2$ theory relevant to the toric code,
\begin{equation*}
\dim \mathcal{H}_{\mathrm{BF},2}(T^2)=4.
\end{equation*}
This is the familiar fourfold topological ground-state degeneracy.

\medskip

\noindent The same answer may also be recognized directly from the discrete gauge-theory viewpoint of Section~\ref{sec:bf-zn}: flat $\Z_N$ gauge fields on $\Sigma_g$ are classified by
\begin{equation*}
\Hom(\pi_1(\Sigma_g),\Z_N)\cong \Z_N^{\,2g},
\end{equation*}
so there are exactly $N^{2g}$ holonomy sectors. Chapter~\ref{ch:dw} will reinterpret that statement as the untwisted Dijkgraaf--Witten partition function on $\Sigma_g\times S^1$, and Chapter~\ref{ch:toporder} will meet the same degeneracy again in the toric-code lattice model.
\end{remark}
